\section{The sample mean as an empirical expectation}

The theoretical expectation \(\E[X]\) is a property of a probability law. It was
introduced as a weighted average of possible values. For observed data, the
corresponding quantity is the sample mean.

\begin{definitionbox}
For a dataset \(x_1,x_2,\ldots,x_n\), the \textbf{sample mean} is
\[
\bar{x}=\frac1n\sum_{i=1}^n x_i.
\]
\end{definitionbox}

The sample mean adds all observed values and divides by the number of
observations. It is a center of the observed data.

\subsection*{A small example}

For the data
\[
2,\quad 4,\quad 4,\quad 7,\quad 9,
\]
the sample mean is
\[
\bar{x}=\frac{2+4+4+7+9}{5}=\frac{26}{5}=5.2.
\]
The number \(5.2\) is not one of the observed values. This is not a problem. The
mean is a balance point of the data, not necessarily a value that appears in the
dataset.

\subsection*{Connection with the empirical distribution}

The sample mean can be interpreted as the expectation with respect to the
empirical distribution. If the empirical distribution assigns mass \(1/n\) to
each observation, then its expectation is
\[
\sum_{i=1}^n x_i\frac1n=\frac1n\sum_{i=1}^n x_i=\bar{x}.
\]
Thus the sample mean is not a new mysterious formula. It is the same idea of
expectation applied to the empirical distribution of the data.

If the data have repeated values, the same idea can be written using frequencies.
Suppose the distinct observed values are \(a_1,\ldots,a_k\), with counts
\(n_{a_1},\ldots,n_{a_k}\). Then
\[
\bar{x}=\sum_{j=1}^k a_j\frac{n_{a_j}}{n}.
\]
This is a weighted average of the distinct observed values, where the weights are
relative frequencies.

\subsection*{Sample mean versus theoretical expectation}

The sample mean \(\bar{x}\) and the theoretical expectation \(\E[X]\) are related
but different.

The expression
\[
\E[X]
\]
belongs to a probability model. It is computed from a probability mass function,
a density, or another description of the law of \(X\). The expression
\[
\bar{x}
\]
belongs to observed data. It is computed from a finite list of numbers.

If the sample is generated from a model with expectation \(\mu=\E[X]\), then
\(\bar{x}\) may be close to \(\mu\), especially for large samples under suitable
assumptions. But for any particular finite sample, \(\bar{x}\) need not equal
\(\mu\).

\begin{remarkbox}
The statement \(\bar{x}=\mu\) is usually false for a finite random sample. The
sample mean is an observed summary; the theoretical mean is a model quantity.
Later chapters explain why sample means often become close to theoretical means
as sample size grows.
\end{remarkbox}

\subsection*{Sensitivity to extreme values}

The sample mean uses all observed values, and each value enters linearly. This
makes it sensitive to extreme observations.

Compare the two datasets
\[
1,\quad 2,\quad 2,\quad 3,\quad 4
\]
and
\[
1,\quad 2,\quad 2,\quad 3,\quad 100.
\]
The first mean is
\[
\frac{1+2+2+3+4}{5}=\frac{12}{5}=2.4.
\]
The second mean is
\[
\frac{1+2+2+3+100}{5}=\frac{108}{5}=21.6.
\]
Changing one observation from \(4\) to \(100\) dramatically changes the mean.
This is not a defect of the arithmetic; it is a property of what the mean
measures. The mean is a balance point, and a large value pulls the balance point
toward itself.

\begin{summarybox}
The sample mean
\[
\bar{x}=\frac1n\sum_{i=1}^n x_i
\]
is the expectation of the empirical distribution. It is the data-based
counterpart of \(\E[X]\), but it is not identical to the theoretical expectation
of a probability model.
\end{summarybox}
